\documentclass[dvipsnames, svgnames, x11names,border=3pt]{standalone}
\usepackage{tikz,pgfplots}
\pgfplotsset{compat = 1.18,width=10cm, height=10cm}


\usetikzlibrary{intersections,math}

\NewDocumentCommand{\iipolt}{ O{name}  m  O{;}}{
\ifdefined\xmin\else
    \def\xmin{\pgfkeysvalueof{/pgfplots/xmin}}
    \def\xmax{\pgfkeysvalueof{/pgfplots/xmax}}
    \def\ymin{\pgfkeysvalueof{/pgfplots/ymin}}
    \def\ymax{\pgfkeysvalueof{/pgfplots/ymax}}
\fi
\addplot+ gnuplot[no markers, raw gnuplot, empty line = jump, name path = #1] {
    set xrange [\xmin : \xmax ];
    set yrange [\ymin : \ymax ];
    set contour;%启用了等高线绘制模式
    unset surface;%禁用了表面绘制模式
    set view map;%将图形视图设置为2D平面视图
    set cntrparam levels discrete 0.0001;%等高线级别设置为0，这样只会绘制函数等于0的等高线，实现了隐函数的绘制
    set isosamples 100;
    set samples 100;
    #2 ;
} #3
}

\begin{document}
\begin{tikzpicture}
\begin{axis}[xmax=24,xmin=0,ymax=5,ymin=0]

\node at(0.5,0.5) {\pgfkeysvalueof{/pgfplots/xmin}}; 

% 隐函数图加 node
\iipolt[a]{splot y - x} [node [pos=0.25,above] {$\sum x_i$}];

% 指定定义域的隐函数图
\iipolt{splot [15:24][0:5]  y - 4* (1 - exp(-(6-y)*(23-x)/18)) +3 } 

% 连续多条隐函数图形
\iipolt{
splot [0:0.2][0:0.4] y - 0.2 - x;
splot [0.4:0.6][0:0.8] y - 0.2 - x
}

% 普通函数图
\iipolt{plot sample [0:0.2] x + 0.4  }

\end{axis}

\end{tikzpicture}
\end{document}

